\documentclass[11pt]{article}
\usepackage[T1]{fontenc}
\usepackage[margin=1in]{geometry}
\usepackage{booktabs}
\usepackage{amsmath}
\usepackage[numbers,sort&compress]{natbib}
\usepackage[hidelinks]{hyperref}
\usepackage{url}
\setlength{\emergencystretch}{3em}
\input{evidence.tex}
\title{PubMed Search MCP: An Auditable Integration Layer for\newline Agent-Assisted Biomedical Literature Retrieval}
\author{Tz-Ping Gau\\Kaohsiung Medical University Hospital}
\date{September 21, 2026\\\small Working manuscript; not submitted}

\begin{document}
\maketitle
\begin{abstract}
Language-model agents can already search the web and call biomedical APIs. The
engineering question is whether a dedicated integration layer improves the
reliability, efficiency, and inspectability of that work. We describe PubMed
Search MCP, a biomedical retrieval service that exposes search, article access,
reference verification, persistent sessions, and reproducible pipelines through
the Model Context Protocol. The service separates host-agent decisions from
retrieval execution and preserves source failures, provenance, and provider
budgets. We distinguish three kinds of existing evidence. A historical lexical
ranking experiment on the BEIR distribution of NFCorpus improved test nDCG@10
from \NdcgBefore{} to \NdcgAfter{}, while the recall@100 interval included zero.
A five-scenario offline scheduling experiment reduced median in-memory MCP
latency by \StaggeredReduction\% in the staggered-dependency case, with unchanged
provider-operation counts and output signatures; other scenarios include a
regression. A three-question native-agent versus packaged-MCP pilot found equal
answer and source-hit scores. These measurements support bounded component and
workflow observations, not an established end-to-end agent advantage. We provide
an evidence-linked manuscript build and a proposed controlled evaluation that
separates server capabilities, client instructions, and provider access.
\end{abstract}

\section{Introduction}
Biomedical literature retrieval combines several distinct activities: translating
a question into queries, locating records, accessing available text, following
references, selecting evidence, and reporting what was searched. A language-model
host can perform these activities using general web search or direct provider
APIs. Adding a domain-specific service therefore requires a comparison with an
already capable baseline, rather than with an agent that has no search tools.

PubMed Search MCP addresses the integration work around retrieval. It provides
consistent article representations, bounded source operations, reusable query
pipelines, and records that a host can inspect after a search. The host remains
responsible for research questions, inclusion criteria, interpretation, and final
claims. The service does not run an autonomous scientific investigation or certify
that a retrieved article supports a clinical conclusion.

This paper makes three engineering contributions: an explicit boundary between
agent reasoning and biomedical retrieval services; inspectable execution and
source-state contracts across search, reading, and export; and a reproducibility
structure that separates component measurements from the unproven causal effect
of the complete agent package. The described software surface is version 0.7.4.
Historical measurements retain their original revision and file fingerprints;
they are not relabeled as measurements of every current capability.

\section{Related work and evaluation targets}
The Model Context Protocol standardizes communication between hosts and tool
servers \citep{mcp2026}. Interoperable transport alone does not determine query
quality, relevance, or appropriate use of evidence. A host can use the same
protocol to access services with substantially different retrieval policies.

The implementation uses established retrieval ideas. BM25 supplies lexical
relevance signals \citep{robertson2009bm25}; reciprocal rank fusion combines ranked
lists \citep{cormack2009rrf}; maximal marginal relevance motivates balancing
relevance and redundancy \citep{carbonell1998mmr}. Their use or combination here is
an implementation choice, not a claim to invent these methods. A scoring option
must also be distinguished from a signal that was actually available and applied
in a particular execution.

Literature search and systematic-review screening have different endpoints.
PRISMA-S specifies reporting considerations for search methods
\citep{rethlefsen2021prismas}; retaining queries and source states can assist that
reporting, but does not establish compliance with the full checklist. ASReview
uses machine learning to prioritize screening within a set of candidate records
\citep{vandeschoot2021asreview}. PubMed Search MCP primarily integrates discovery
and access, and cannot infer systematic-review completeness from a short ranked
result list.

Public benchmarks test different units. BEIR supports retrieval comparisons
across datasets \citep{thakur2021beir}; NFCorpus supplies a medical retrieval task
\citep{nfcorpus}. LitSearch contains 597 scientific literature queries in machine
learning and natural language processing \citep{ajith2024litsearch}, making domain
transfer an explicit concern for a biomedical service. ScholarGym separates query
planning, tool invocation, and relevance assessment in a static environment with
approximately 570,000 papers and 2,536 queries \citep{shen2026scholargym}.
PaperSearchQA studies scientific search and question answering
\citep{burgess2026papersearchqa}. An adapted tool interface or corpus changes the
evaluation setting and must be disclosed rather than presented as an official
leaderboard result.

OpenResearch offers a broader workspace for agent-driven research and experimental
artifacts \citep{openresearch2026}. Its emphasis on inspectable artifacts is useful
context for integration design. Its scope differs from a retrieval service, and
we report no head-to-head quality comparison with it.

\section{Service design}
\subsection{Host and service responsibilities}
Version 0.7.4 exposes 41 MCP tools organized into 16 registry categories through
the MCP SDK v2 \texttt{MCPServer}. Presentation wrappers validate requests and
delegate to application services. Domain models represent articles, pipelines,
and research artifacts; infrastructure adapters implement external providers
and storage. Keeping this direction of dependency allows transport checks and
retrieval checks to exercise separate contracts.

The host supplies a query or a pipeline specification. For a PICO workflow, the
host supplies the population, intervention, comparison, and outcome fields; the
backend does not silently infer those fields from an unrestricted question.
Likewise, host instructions and optional skills can guide tool use, but are
separate from the server implementation. Their contribution must be isolated in
an evaluation intended to explain an observed package-level gain.

\begin{table}[htbp]
\centering\small
\begin{tabular}{p{0.23\linewidth}p{0.68\linewidth}}
\toprule
Boundary & Responsibility and observable output \\
\midrule
Host agent & Plans queries, chooses articles, interprets findings, writes claims. \\
MCP presentation & Validates parameters, exposes the registry, serializes results. \\
Application services & Execute searches and pipelines, aggregate articles, manage sessions and exports. \\
Domain and infrastructure & Represent evidence objects; access providers and storage under provider-specific constraints. \\
\bottomrule
\end{tabular}
\caption{Responsibility boundaries. A successful tool call does not validate the host's scientific conclusion.}
\end{table}

\subsection{Retrieval, access, and provenance}
Multi-source retrieval must reconcile identifiers and metadata without turning
provider-specific absence into a universal assertion. A result may lack an
abstract, a citation metric, or accessible full text. Source diagnostics separate
an empty response from an unavailable, failed, or unrequested source. Ranking
provenance distinguishes the requested configuration from effective signals.
For example, a citation-based ordering should not be reported as applied when
the metric provider returned no usable values.

Article access and reference verification remain separate operations. Resolving
an identifier or matching bibliographic fields checks identity and consistency;
it does not establish support for a statement. Claim validation requires access
to the relevant passage, its context, and a judgment about the claim. Similarly,
exporting a note or a CSL record preserves metadata rather than completing
scientific verification. Unverified status and the option to omit abstracts must
survive export.

Persistent sessions and pipeline artifacts let a host revisit search decisions
without reconstructing them from conversational text alone. Research-chronicle
views organize articles and relationships over time. Their importance signals,
milestones, and inferred relationships are descriptive heuristics: they are not
validated measures of research maturity, causal influence, or clinical evidence
strength. A citation graph also does not by itself identify support, contradiction,
or the reason one article cites another.

\subsection{Dependency scheduling and provider protection}
A pipeline describes steps and dependencies. A dependency-ready scheduler can
start an eligible child after its prerequisites finish, without waiting for an
unrelated slow branch at the same nominal level. The useful optimization is less
avoidable waiting, not unrestricted growth in request concurrency.

Provider protection combines applicable rate admission, bounded in-flight
operations, shared cooldown behavior, retry handling, and cancellation cleanup.
The documented scope matters: limits coordinated within a process or event loop
do not constitute a distributed quota across independently deployed workers.
An operator sharing one provider account across processes needs an appropriate
shared admission mechanism or a deployment-wide conservative allocation.
No result in this paper demonstrates permission to exceed an upstream quota.

Transport and artifact ownership also affect safe integration. Local and service
deployments have different authentication and filesystem boundaries. Client
harness installation preserves existing user-owned instruction files and whole
skill directories by default. Repository updates to canonical instructions are
not authorization to replace a user's customized installed copy.

\section{Evidence and methods}
We use four existing evidence records: lexical retrieval, offline execution,
a small agent pilot, and release validation. An evidence manifest pins their
SHA-256 hashes. The manuscript builder reads numerical tables from these records
and rejects modified inputs. This protects against transcription drift; it does
not independently validate the original experiments or eliminate bias.

\subsection{Historical lexical retrieval experiment}
The September 9 experiment evaluated the production lexical scorer against a
baseline revision beginning \texttt{dbcf0c8}. It used the BEIR distribution of
NFCorpus: 3,633 documents, 324 development queries, and 323 test queries. Full-corpus
document frequencies and average length were available in both conditions, with
deterministic identifier tie-breaking. The recorded protocol uses linear-gain
nDCG and treats unjudged documents as zero gain. Recall is consequently recall
against the available judgments, not an exhaustive measure of all relevant
biomedical literature.

The experiment did not invoke an LLM, external search APIs, full-text reading,
or the complete MCP package. Its paired bootstrap intervals use 10,000 resamples
and seed 20260909. Original data hashes, per-query values, baseline commit, and
before/after production-file hashes remain in the report. This is historical
component evidence, not a fresh evaluation of all version 0.7.4 behavior.

\subsection{Offline execution experiment}
The September 18 execution benchmark used the real pipeline executor and an
in-memory MCP client, replacing providers with fixed-delay fixtures. Twenty
repetitions per condition covered no-I/O, single-step, balanced-branch,
staggered-branch, and serial-dependency scenarios. The fixture records provider
operations, successful steps, output identifiers, and latency samples. Comparing
signatures checks that a timing change did not omit required work.

This setup excludes live provider latency, host-model planning, network transport,
process startup, and real rate-limit incidents. Session history grows within the
fixture, so the difference between executor and MCP timings is not a clean estimate
of protocol overhead alone. The result records include executor and script hashes;
the shared recorded Git revision alone does not identify the uncommitted change
that was measured.

\subsection{Exploratory agent pilot and release checks}
The September 9 agent pilot selected three queries from the 5,000-query
PaperSearchQA test split, with one run per condition. Both conditions used
\texttt{gpt-6-astra}, xhigh reasoning effort, and native web search; the package
condition additionally had the historical MCP package. The corpus was not frozen
and token and backend-call budgets were not enforced. This pilot checks that
comparative execution and scoring are possible, with substantial confounding and
insufficient sample size for a general performance claim.

Release validation provides a different kind of evidence. The version 0.7.4
receipt records 4,758 passing local tests, 23 skips, and 30 deselections under
Python 3.10 and 3.13. It also records actual stdio, Streamable HTTP, and installed-wheel
checks of the 41-tool surface. These are software-contract checks, not measures
of scientific relevance. Counts alone do not measure test value; maintaining
public contracts and concrete regression coverage is the objective.

\section{Results}
\subsection{Lexical relevance}
Table~\ref{tab:retrieval} reports the historical lexical comparison. Test nDCG@10
increased by \NdcgDelta{} (paired 95\% bootstrap interval
[\NdcgLow{}, \NdcgHigh{}]). Of 323 queries, 49 improved, 33 regressed, and 241 were
unchanged. The absolute change is small and specific to this corpus and scorer.
Test recall@100 changed by \RecallDelta{}, with interval
[\RecallLow{}, \RecallHigh{}], which includes zero. We therefore do not infer a
recall improvement, clinical coverage improvement, or agent-level benefit from
this experiment.
\input{retrieval-table.tex}

\subsection{Execution latency}
Table~\ref{tab:execution} includes every measured scenario. In the staggered case,
the median executor time fell from 202.067 to 121.889 ms, and median MCP time fell
from 391.957 to 314.257 ms (\StaggeredReduction\% reduction). Both conditions
performed four provider operations and returned the same output signature. The
serial case became slower at the MCP level; performance is not uniformly improved.
These observations support an implementation-level reduction in waiting for this
fixture, without predicting the same percentage for a live literature search.
\input{execution-table.tex}

\subsection{Agent-level uncertainty}
Both pilot conditions answered two of three queries exactly and retrieved the
gold source PMID for one of three. The observed paired differences were zero.
With three pairs, neither equality of the systems nor a positive MCP effect is
established. We omit the degenerate bootstrap intervals stored for this pilot
because they would invite unjustified precision. No complete 5,000-query run,
current-package comparison, or official benchmark submission is reported.

\section{Proposed controlled evaluation}
The key estimand is the difference between a capable native host and that same
host with the integration package under comparable information and resource
access. A proposed three-arm evaluation uses (A) native search and direct PubMed
access, (B) the original MCP package, and (C) the current package. The model,
reasoning configuration, task text, and availability of native search should be
matched. Arm B must be identified by an immutable revision, and C must be built
from a pinned revision rather than a changing working tree.

Two tracks answer different questions. A frozen-corpus track supplies all arms
with the same article and reference snapshot through instrumented adapters. It
estimates the benefit of planning support, aggregation, state, and tool design
conditional on that snapshot. An ecological live-search track retains actual
web and provider access and measures the whole deployment experience, including
coverage differences. The tracks must be reported separately; replacing a native
search backend with a fixture changes the evaluated product.

The main retrieval endpoints should be preregistered for the selected task:
for ranked retrieval, nDCG@10 and judged recall@100; for high-recall screening,
recall and work saved at a prespecified sensitivity; for question answering,
answer and supporting-source scores. A claim-support subset requires passages
and an adjudication rubric that distinguishes support, contradiction, insufficient
context, and inaccessible evidence. Bibliographic matches cannot substitute for
that annotation.

Resource accounting must occur below the MCP interface. One bundled tool call
may issue several searches, reads, or retries. Record provider requests, unique
documents read, returned text volume, model tokens, latency, failures, and cache
hits. Pair questions across arms, isolate sessions and caches, randomize execution
order, and include failed or timed-out runs in the denominator. Analyze paired
differences with uncertainty estimates and disclose any development-set tuning.
An additional server-only versus server-plus-instructions ablation separates
service utility from prompt guidance.

All executions must respect provider-specific admission and cooldown policies.
A repeatable benchmark should use frozen fixtures for broad sweeps and a small,
separately budgeted live validation for provider contracts. The proposed protocol
is documented in the repository but has not been preregistered or fully executed.
This paper does not claim results for that future experiment.

\section{Limitations and reproducibility}
The existing evidence is heterogeneous in revision, scope, and environment. The
retrieval judgments are incomplete; the timing fixture is synthetic; the pilot
is small and uses live search. None establishes improved systematic-review recall,
reliable citation sentiment, or more accurate scientific conclusions. Provider
availability and model behavior can change after a release. Loop-local admission
does not guarantee safety across a distributed fleet, and local software tests
cannot prove that an external service will never throttle a request.

The repository publishes software under Apache-2.0 and records its release in
\texttt{CITATION.cff}. The current manuscript, reviewed bibliography, evidence
manifest, and build script are maintained in one publication workspace. Builds
generate tables from pinned evidence and produce a self-contained TeX source
archive with the resolved bibliography. Historical drafts are archived separately
and are not evidence for present claims. The manuscript has no assigned arXiv
identifier or publication DOI at the time of writing.

Software and evidence are available at
\url{https://github.com/u9401066/pubmed-search-mcp}. Relevant reproducibility
entry points are \path{scripts/build_publication.py},
\path{docs/publication/evidence.json}, and the dated records in
\path{docs/reports/}. Build receipts identify the source and input hashes;
they do not imply that a later commit reproduces historical experiments unchanged.

\section{Conclusion}
PubMed Search MCP packages biomedical retrieval as an inspectable service for a
host agent. Its contribution is integration: explicit responsibility boundaries,
source-state reporting, reusable execution, and bounded provider operations.
Existing records show a small lexical-ranking change and a scenario-dependent
scheduling improvement, while the agent pilot establishes no overall advantage.
A controlled, resource-accounted comparison is still needed to quantify what the
complete package adds over native search.

\bibliographystyle{plainnat}
\bibliography{references}
\end{document}
